Source : DGFiP, Statistiques n° 51 (septembre 2026), tableau 2, exercice 2025 (données provisoires).\par Note : IS brut, avant crédits et réductions d'impôt (70,5 Md\euro{} en 2025). Catégories d'entreprises du décret n° 2008-1354 (une entreprise peut regrouper plusieurs sociétés). Le résultat fiscal (bénéfices moins déficits) se répartit en 17, 26, 24 et 34 \%. En 2022, la répartition de l'IS brut était presque la même (19, 28, 23 et 30 \%). La contribution exceptionnelle de 2025 et 2026 (7,5 Md\euro{} en 2025) s'ajoute pour les seuls groupes de plus de 1 Md\euro{} de chiffre d'affaires.
